%! Parametric Functions and Rates of Change — Unit 4, Lesson 4.3 — Group Activity
\documentclass[10pt]{article}
\usepackage{precalculus-article}
\usepackage{precalculus-boxes}
\newcommand{\TallMath}[1]{$\displaystyle #1\rule[-1.4em]{0pt}{3.2em}$}

\begin{document}

\pageheader{Unit 4, Lesson 4.3}{Group Activity}
\namepartnerperiod

\vspace{0.12in}

\begin{scenariobox}[Two Robots, One Path]{sky}
Two robots traverse the same parabolic path in a warehouse.
\[
  \text{Robot A:}\quad g_A(t) = (t,\; t^2), \quad 0 \le t \le 3
  \qquad\qquad
  \text{Robot B:}\quad g_B(t) = (3-t,\; (3-t)^2), \quad 0 \le t \le 3
\]
Robot A starts at $(0, 0)$ and Robot B starts at $(3, 9)$.
Your goal is to analyze their directions, the path they share, and
the slope of the curve between two points.
\end{scenariobox}

\vspace{0.10in}

\begin{notesbox}{Part 1 --- Tables and Shared Path}

\textbf{1.}\quad Complete the table for each robot. Leave no cell blank.

\vspace{6pt}
\begin{center}
\renewcommand{\arraystretch}{1.7}
\begin{tabular}{|c|c|c|}
\hline
$t$ & Robot A position $(x, y)$ for $g_A(t) = (t,\; t^2)$ &
      Robot B position $(x, y)$ for $g_B(t) = (3-t,\; (3-t)^2)$ \\
\hline
$0$ & \blank{3cm} & \blank{3cm} \\
\hline
$1$ & \blank{3cm} & \blank{3cm} \\
\hline
$2$ & \blank{3cm} & \blank{3cm} \\
\hline
$3$ & \blank{3cm} & \blank{3cm} \\
\hline
\end{tabular}
\end{center}

\vspace{8pt}
\textbf{2.}\quad Do both robots visit the same set of ordered pairs?
Explain using your table.

\writelines{2}

\end{notesbox}

\vspace{0.08in}

\begin{notesbox}{Part 2 --- Direction of Motion}

\textbf{3.}\quad At $t = 1$, is Robot A moving right or left? Up or down?
Show your reasoning by comparing values in the table.

\writelines{3}

\textbf{4.}\quad At $t = 1$, is Robot B moving right or left? Up or down?
Show your reasoning by comparing values in the table.

\writelines{3}

\textbf{5.}\quad Do the two robots move in the same direction at $t = 1$?
Explain what this tells you about the two parametrizations.

\writelines{2}

\end{notesbox}

\vspace{0.08in}

\begin{notesbox}{Part 3 --- Slope of the Curve}

\textbf{6.}\quad Using Robot A's data, compute the slope of the curve
between $t = 0$ and $t = 2$.

\vspace{4pt}
$x_A(0) = $ \blank{1.5cm} \qquad $x_A(2) = $ \blank{1.5cm} \qquad $\Delta x = $ \blank{1.5cm}

$y_A(0) = $ \blank{1.5cm} \qquad $y_A(2) = $ \blank{1.5cm} \qquad $\Delta y = $ \blank{1.5cm}

\vspace{4pt}
Slope $= \dfrac{\Delta y}{\Delta x} = $ \blank{2.5cm}

\vspace{10pt}
\textbf{7.}\quad Using Robot B's data, compute the slope of the curve
between the corresponding two points (the same two $(x,y)$ locations).
Do you get the same slope? Explain why this makes sense.

\vspace{4pt}
Robot B passes through $(0, 0)$ at $t = $ \blank{1.2cm}
and through $(2, 4)$ at $t = $ \blank{1.2cm}.

$\Delta x = $ \blank{1.5cm} \qquad $\Delta y = $ \blank{1.5cm}

\vspace{4pt}
Slope $= \dfrac{\Delta y}{\Delta x} = $ \blank{2.5cm}

\writelines{2}

\end{notesbox}

\vspace{0.08in}

\begin{notesbox}{Part 4 --- When Does Each Robot Arrive?}

\textbf{8.}\quad At what value of $t$ does Robot A pass through $(1, 1)$?
At what value of $t$ does Robot B pass through $(1, 1)$?
Show your work (set the component equations equal to $1$ and solve).

\writelines{4}

\textbf{9.}\quad At what value of $t$ does each robot pass through
$(0, 0)$? Through $(3, 9)$?

\writelines{3}

\end{notesbox}

\end{document}
